\exheader

\begin{exercise}
Decide if each table could represent a linear function and explain your reasoning.  For each table that could represent a linear function, find a formula that matches the data given in the table.

\begin{parts}
    \item
\begin{tabular}{|c|c|c|c|c|}
\hline $x$&0&1&2&3
\\\hline $y$&12.6&9.3&6.0&2.7
\\\hline
\end{tabular}

\item \begin{tabular}{|c|c|c|c|c|c|cc c|}
\hline $t$ & $-1$ & 1 & 3 & 5 
\\\hline $f(t)$ & 0.1 & 0.4 & 0.7 & 1.0 
\\\hline
\end{tabular}

\item \begin{tabular}{|c|c|c|c|c|c|c|cc c|}
\hline $n$ & 0 & 1.5 & 3 & 4.5 & 6 
\\\hline $g(n)$ & 12 & 18 & 12& 6& 12
\\\hline
\end{tabular}

\item \begin{tabular}{|c|c|c|c|c|c|c|cc c|}
\hline $t$ & 0 & 2 & 4 & 6 & 8 
\\\hline $h(t)$ & 6.4 & 1.6 & 0.4& 0.1& 0.025
\\\hline
\end{tabular}

\item \begin{tabular}{|c|c|c|c|c|c|cc c|}
\hline $z$ & $-2$ & 1 & 4 & 7 
\\\hline $w(z)$ & 6 & 4.8 & 3.6 & 2.4 
\\\hline
\end{tabular}
\end{parts}

\end{exercise}

\begin{solution}
    \begin{parts}
        \item linear with equation $y=-3.3x+12.6$
        \item linear with equation $f(t)=0.15t+0.25$
        \item not linear
        \item not linear
        \item linear with equation $w(z)=-0.4z+5.2$
    \end{parts}
\end{solution}



\begin{exercise}
The table below shows for each temperature $F$ in degrees Fahrenheit the corresponding approximate temperature $C$ in degrees Celsius. Is it appropriate to model $C$ by a linear function of $F$? If it is, find a formula for $C$ as a function of $F$.
\begin{center}
\begin{tabular}{|c|c|c|c|c|c|cc c|}
\hline $F$ (degrees Fahrenheit) & 32 & 42 & 52 & 62 
\\\hline $C$ (degrees Celsius) & 0 & 5.55 & 11.12 & 16.67 
\\\hline
\end{tabular}
\end{center}
\end{exercise}

\begin{solution}
approximately linear with approximate equation $C=0.555F-17.76$
\end{solution}



\begin{exercise}
The height of a human individual can be estimated by the length of the femur, as shown for males in the following table.
\begin{center}
\begin{tabular}{|c|c|c|c|c|}
\hline $L$ (length of femur in cm) & 42 & 45 & 48  & 51
\\\hline $H$ (height of human male in cm) & 162.97 & 169.93 & 176.89 &183.85
\\\hline
\end{tabular}
\end{center}

\begin{parts}
    \item Use the information in this table to find a possible formula for the height of a human male $H$ as a function of his femur length $L$.
    
    \item What would the approximate height of a human male with femur length 46 cm be?
    
    \item Fernando is 175.4 cm tall. What would you expect the approximate length of his femur to be?
\end{parts}
\end{exercise}

\begin{solution}
\begin{parts}
    \item $H=2.32L+65.53$
    \item $172.25$ cm
    \item approximately 47.36 cm
\end{parts}
\end{solution}




\begin{exercise}
    The weight of water above a scuba diver as well as the air above the diver exerts pressure on their bodies. The pressure the diver experiences at sea level is 14.7 PSI (pounds per square inch), and this pressure increases by 0.4 PSI per each foot of depth.

    \begin{parts}
        \item  Write a linear equation expressing the pressure $P$ on a diver at a depth of $d$ feet below sea level.

        \item The deepest a recreational scuba diver typically dives is 130 feet. What is the pressure on a diver at this depth?
    \end{parts}
\end{exercise}


\begin{solution}
    \begin{parts}
        \item $P=0.4d+14.7$
        \item 66.7 PSI
    \end{parts}
\end{solution}




\begin{exercise}
    A hiker is at a trailhead about to climb a mountain. The temperature at the trailhead is $65^{\circ}$F. According to the standard atmosphere model\footnote{https://www.grc.nasa.gov/WWW/K-12/airplane/atmos.html, accessed: 6/25/20}, the temperature drops by $0.00356^{\circ}$F per each 1 foot increase in altitude. Let $T=T(h)$  be the temperature, in $^{\circ}$F, $h$ feet above the hiker.

    \begin{parts}
        \item Write a formula for the function $T(h)$.

         \item Find the temperature at the mountaintop that is $6210$ feet above the hiker.

        \item Is the function $T(h)$ increasing or decreasing?
    \end{parts}

\end{exercise}

\begin{solution}
    \begin{parts}
        \item $T(h)=-0.00356h+65$
        \item $42.89$ $^\circ$ F
        \item The function $T(h)$ is decreasing.
    \end{parts}
\end{solution}




\begin{exercise}
    A man goes to a gym to exercise. After $t$ minutes on a treadmill, his pulse (heart rate), $H$, in beats per minute, is: 

\begin{center}
\begin{tabular}{|c|c|c|c|c|c|c|}
\hline $t$ (minutes) & 0 & 3 & 6 & 9 & 12 & 15 
\\\hline $H$ (bpm) & 85 & 87 & 89 & 91 & 93 & 95
\\\hline
\end{tabular}
\end{center}

\begin{parts}
    \item Is it appropriate to model $H(t)$ by a linear function? If yes, find a formula for $H(t)$.
    
    \item If the function $H(t)$ can be modeled by a linear function, give units for the slope and the vertical intercept of $H(t)$. Then explain their meaning in practical terms.
\end{parts}
\end{exercise}

\begin{solution}
    \begin{parts}
        \item Linear; formula $H(t)=\frac{2}{3}t+85$

        \item The slope is $\frac{2}{3}$ bpm per minute; the practical meaning of the slope is that the man's bpm is increasing at a rate of $\frac{2}{3}$ bpm per minute. The vertical intercept is $85$ bpm; the practical meaning of the vertical intercept is that the man's resting heart rate (heart rate when performing no exercise on the treadmill) is $85$ bpm.
    \end{parts}
\end{solution}